\documentclass[10pt]{article}
\usepackage[margin=0.78in]{geometry}
\usepackage{fontspec}
\defaultfontfeatures{Ligatures=TeX}
\setmainfont{Liberation Serif}
\setsansfont{DejaVu Sans}
\setmonofont{DejaVu Sans Mono}
\usepackage{amsmath,amssymb}
\usepackage{hyperref}
\setlength{\parindent}{0pt}
\setlength{\parskip}{0.5em}
\begin{document}

\begin{center}
{\Large \textbf{Theory Report}}\\[0.35em]
{\large Derivation of Gibbs measure from Gibbs state with the fractional Bessel interaction in Two Dimensions}\\[0.45em]
Phan Thanh Nam, Rongchan Zhu, and Xiangchan Zhu\\
\href{https://arxiv.org/abs/2604.21583}{arXiv:2604.21583} \hspace{1em} Digest date: 2026-04-24
\end{center}

\section*{Paper information}
This paper studies a quantum-to-classical limit for a two-dimensional bosonic grand-canonical Gibbs state on the torus $\mathbb T^2$. The interaction is the fractional Bessel potential whose Fourier coefficients are
\[
\widehat v_\beta(k)=\langle k\rangle^{-\beta}=(1+|k|^2)^{-\beta/2},\qquad \tfrac32<\beta\le 2,
\]
so $v_\beta=(1-\Delta)^{-\beta/2}$ as an operator kernel on $L^2(\mathbb T^2)$. The range $\beta\le 2$ is exactly the non-summable regime in dimension two, so
\[
\sum_{k\in\mathbb Z^2}\widehat v_\beta(k)=\infty,
\]
which means the naive density-square rewriting of the quantum interaction produces an infinite self-energy.

The quantum model lives on bosonic Fock space with one-body Hamiltonian $h=-\Delta+1$. At temperature of order $\lambda^{-1}$, the authors consider a renormalized Gibbs state
\[
\Gamma^{\mathrm{re}}_\lambda=\frac{e^{-H^{\mathrm{re}}_\lambda}}{Z^{\mathrm{re}}_\lambda},
\qquad
H^{\mathrm{re}}_\lambda=\lambda d\Gamma(h)+W^{(2)}_\beta+c_\lambda N+C_\lambda,
\]
where the renormalization is implemented by extracting only the zero Fourier mode and replacing it by a centered number-fluctuation term proportional to $(N-N_0)^2$. The target object is the classical Gibbs measure
\[
 d\nu(u)=z_{\mathrm{cl}}^{-1}e^{-\mathcal D(u)}\,d\mu_0(u),
\]
where $\mu_0$ is the Gaussian free field with covariance $h^{-1}$ and $\mathcal D(u)$ is the Wick-renormalized fractional-Bessel Hartree energy.

\section*{Executive summary}
The main theorem says that as $\lambda\downarrow 0$, the renormalized quantum Gibbs state converges to the classical Gibbs measure associated with the same interaction, simultaneously at the level of free energy and correlation operators. More concretely,
\[
-\log\frac{Z^{\mathrm{re}}_\lambda}{Z_0}\to -\log z_{\mathrm{cl}},
\qquad
k!\lambda^k\Gamma^{(k)}_\lambda\to \int |u^{\otimes k}\rangle\langle u^{\otimes k}|\,d\nu(u)
\]
in Hilbert--Schmidt norm for each fixed $k\ge 1$.

What makes the result nontrivial is not merely that the interaction is nonlocal, but that it is singular enough that the quantum Hamiltonian cannot be put into the standard positive density-square form. The obstruction is the divergent self-energy coming from the Fourier sum of $\widehat v_\beta$. So the paper has to do two hard things at once: define the correct renormalized many-body model, and show that its variational structure still matches the classical fractional-Bessel Gibbs theory in the weak-coupling/high-temperature limit.

The proof architecture is very clean. First, localize to low Fourier modes. Second, compare the low-mode quantum problem with a finite-dimensional classical Hartree problem through coherent states, lower symbols, Berezin--Lieb, and a quantitative quantum de Finetti theorem. Third, prove that every high-frequency remainder produced by localization vanishes. The paper's real contribution is that third step: it develops a multi-scale decomposition of the ultraviolet error into zero-mode, shell, and far-tail pieces and controls each by a different mechanism.

\section*{Problem setup and intuition}
The classical interaction functional is formally
\[
\mathcal D(u)=\frac12\int_{\mathbb T^2\times\mathbb T^2} v_\beta(x-y):|u(x)|^2::|u(y)|^2:\,dx\,dy,
\]
with Wick ordering relative to the Gaussian free field. Because $u$ under $\mu_0$ is distribution-valued, this renormalization is essential even classically.

On the quantum side, the genuine two-body operator is harmless in quartic form but dangerous in density form. If one writes the density Fourier modes
\[
\rho_k=d\Gamma(e_k)=\frac1{2\pi}\sum_{p\in\mathbb Z^2} a^*_{p+k}a_p,
\]
then the canonical commutation relations give
\[
W^{(2)}_\beta \sim \frac{\lambda^2}{2}\sum_k \widehat v_\beta(k)\rho_k\rho_{-k}
-\frac{\lambda^2}{2(2\pi)^2}\Big(\sum_k \widehat v_\beta(k)\Big)N.
\]
The second term diverges because $\widehat v_\beta\notin \ell^1(\mathbb Z^2)$. So one cannot start from the usual positive density-square representation. The authors instead keep all nonzero modes in the genuine quartic operator and renormalize only the zero mode by replacing the raw quadratic number term with a centered fluctuation $(N-N_0)^2$.

The intuition behind the limit is standard in the nonlinear Gibbs program: when temperature is of order $\lambda^{-1}$ and the coupling is scaled as in the Hamiltonian, the bosonic Gibbs state should behave semiclassically and be described by a random field rather than a many-body density operator. But here that transition has to survive an ultraviolet singularity that is still present at the quantum level. The proof therefore asks: after projecting to low frequencies, can the leftover ultraviolet pieces be shown to disappear fast enough to recover the classical functional exactly?

\section*{Main theorem/result (informal but precise)}
For every fixed $\beta\in(\tfrac32,2]$, the authors prove two coupled convergence statements.

\textbf{1. Relative free energy convergence.} Let $Z^{\mathrm{re}}_\lambda$ be the partition function of the renormalized interacting Bose gas and $Z_0$ that of the free Gibbs state. Then
\[
-\log(Z^{\mathrm{re}}_\lambda/Z_0)\to -\log z_{\mathrm{cl}}.
\]
So the quantum variational problem converges to the classical fractional-Bessel free-energy problem.

\textbf{2. Density-matrix convergence.} For every fixed $k\ge 1$,
\[
k!\lambda^k\Gamma^{(k)}_\lambda \to \gamma^{(k)}_{\nu}:=\int |u^{\otimes k}\rangle\langle u^{\otimes k}|\,d\nu(u)
\]
in Hilbert--Schmidt norm on the symmetric $k$-particle space. In other words, the rescaled quantum $k$-point functions converge to the $k$-th moments of the limiting classical field.

A key point is what is actually rigorous. The theorem is proved for the full non-summable range $\tfrac32<\beta\le 2$. The restriction $\beta>\tfrac32$ is not cosmetic: it enters the shell summability estimates and the control of the far-tail block. The authors explicitly note that new ideas would be needed below that threshold.

\section*{Proof architecture}
\textbf{Step 1: Variational setup and a priori bounds.} The argument begins from the Gibbs variational principle. Even the preliminary bounds are Bessel-specific because one cannot harmlessly rewrite the interaction as a density square. The nonzero modes are estimated directly from the quartic formula, and the cancellation $\mathrm{Tr}(\gamma_0 M_k)=0$ for $k\neq 0$ removes the dangerous direct Wick term in source-perturbed estimates.

\textbf{Step 2: Low-frequency localization.} Choose a low-frequency projector $P=\mathbf 1_{\{h\le \Lambda^2\}}$ with $\Lambda=\lambda^{-\alpha}$, $0<\alpha<1/4$. The localized state on $P\mathcal H$ is compared with a finite-dimensional Hartree functional
\[
\mathcal D_P(u)=\frac12\sum_{k\ne 0}\widehat v_\beta(k)|\langle u,e_{k,P}u\rangle|^2
+\frac1{2(2\pi)^2}(\|u\|^2-(2\pi)^2\rho_{0,P})^2.
\]
This is the classical effective energy seen by the low modes.

\textbf{Step 3: Quantum-to-classical comparison on the low modes.} The authors pass from the localized quantum state to a probability measure on $P\mathcal H$ using coherent states and the lower symbol (Husimi measure). The quantitative de Finetti theorem expresses low-order reduced density matrices in terms of moments of that lower symbol, with explicit error terms. Berezin--Lieb and entropy monotonicity then give a free-energy lower bound by the classical variational functional. A pointwise comparison between the free lower symbol and the Gaussian free field measure shows that the free reference state really becomes the classical Gaussian measure on the projected space.

\textbf{Step 4: Decompose the ultraviolet error.} Localization produces explicit high-frequency remainders. These are split into three pieces: (i) a zero-mode high-frequency remainder, (ii) a shell contribution after introducing a second cutoff $\Lambda_1=\lambda^{-\alpha_1}$ with $1/2<\alpha_1<3/4$, and (iii) a far-tail contribution beyond that shell.

\textbf{Step 5: Control the zero-mode remainder.} This is handled by a second-order correlation inequality plus a careful block decomposition of the relevant double commutator. The useful structural fact is that every surviving quartic monomial still contains a low-frequency leg, which lets the authors estimate it by asymmetric quartic Hilbert--Schmidt bounds together with kinetic-energy estimates.

\textbf{Step 6: Show the shell contribution vanishes.} The shell piece is not positive and not termwise negligible. The authors rewrite it in terms of shell-mode fluctuation observables, polarize those observables, and apply the same commutator strategy mode by mode. This is where the restriction $\beta>3/2$ appears in a weighted summability estimate.

\textbf{Step 7: Control the far tail from below.} The far-tail part is treated differently. After discarding a positive high-high block, the mixed remainder is kept as a genuine off-diagonal two-body operator. Its operator norm is paired with the small high-frequency particle-number mass, yielding a quantitative lower bound $E_{\mathrm{tail}}(\Gamma_\lambda)\ge -o(1)$.

\textbf{Step 8: Match upper and lower bounds.} The free-energy lower bound comes from the classical variational principle after all remainders are shown negligible. For the upper bound, the authors build a trial state that is interacting on the low-frequency Fock space and free on the high-frequency complement. Because this ansatz is compatible with localization, the same error decomposition works and the remainder terms vanish again. The projected partition function converges to the finite-dimensional classical partition function, which then converges to $z_{\mathrm{cl}}$.

\textbf{Step 9: Recover density matrices.} Once the free energy converges, the lower symbols converge to the projected classical Gibbs measures with polynomial weights. The quantitative de Finetti formula converts this weighted measure convergence into scalar moment convergence, and then into Hilbert--Schmidt convergence of the full rescaled reduced density matrices.

\section*{Why it matters / limitations / takeaway}
This paper pushes the nonlinear Gibbs derivation program into a regime where the quantum interaction is itself singular enough to break the standard renormalization template. That is the main conceptual achievement. It shows that the bridge from quantum bosonic Gibbs states to classical field theory does not depend on having a summable interaction kernel or an immediately positive density-square representation.

The main limitation is also clear. The proof needs $\beta>3/2$, and the authors explain that below this threshold both shell and tail estimates would need genuinely new ideas. The setting is also homogeneous on the torus; extensions to inhomogeneous traps or to the true Coulomb interaction are discussed as plausible but not carried out here.

The clean takeaway is: after the correct zero-mode renormalization, the two-dimensional Bose gas with fractional-Bessel interaction has the expected classical fractional-Bessel Gibbs measure as its precise weak-coupling/high-temperature limit. The paper proves this not by a black-box semiclassical argument, but by a sharp variational reduction plus a delicate ultraviolet remainder analysis. That remainder analysis is the real engine of the theorem.

\end{document}
